%% ----------------------------------------------------------------------------
%%  §1.12  Punti fondamentali del capitolo
%% ----------------------------------------------------------------------------
\section{Punti fondamentali del capitolo}
\label{sec:sintesi-cap01}

Il capitolo ha fissato le regole del gioco: gli undici assiomi dei numeri reali, tra i
quali l'assioma di completezza è l'unico che distingue $\R$ da $\Q$. Su questa base
abbiamo costruito il linguaggio che useremo in tutto il libro (insiemi, funzioni,
induzione) e il catalogo delle funzioni elementari.

\begin{puntichiave}
\item \textbf{Gli assiomi.} $\R$ soddisfa le proprietà di campo (A1)--(A6), le proprietà
dell'ordinamento (O1)--(O4) e l'assioma di completezza (\cref{ass:completezza}): se $A$ e
$B$ sono non vuoti e $a\le b$ per ogni $a\in A$ e $b\in B$, esiste $c\in\R$ con
$a\le c\le b$ per ogni $a\in A$ e $b\in B$. Ogni altra proprietà si dimostra.
\item \textbf{Le regole dell'algebra sono teoremi.} La legge di annullamento del prodotto
(\cref{prop:annullamento}), la regola dei segni (\cref{prop:segni}) e il cambio di verso
di una disuguaglianza moltiplicata per $c\le0$ (\cref{prop:ordine-prodotto}) discendono
dagli assiomi; la divisione per zero è impossibile perché porterebbe a $0=1$.
\item \textbf{Il linguaggio degli insiemi.} Due insiemi si dimostrano uguali per doppia
inclusione; le classi di una relazione di equivalenza formano una partizione
(\cref{prop:classi}); la relazione $\le$ è un ordine totale su $\R$.
\item \textbf{$\Q$ non è completo.} Nessun razionale ha quadrato $2$
(\cref{teo:radice2}), e gli insiemi \eqref{eq:AB-radice2} non sono separati da alcun
razionale (\cref{teo:Q-non-completo}): $\Q$ soddisfa (A1)--(A6) e (O1)--(O4), ma non la
completezza, che in $\R$ garantisce l'esistenza di $\sqrt2$.
\item \textbf{Funzioni.} Una funzione associa a ogni elemento del dominio \emph{uno e un
solo} valore; è invertibile se e solo se è biunivoca, e allora
$f^{-1}\circ f=\mathrm{id}_{A}$ e $f\circ f^{-1}=\mathrm{id}_{B}$.
\item \textbf{Criterio di invertibilità.} Una funzione $f\colon A\to B$ suriettiva e
strettamente monotòna è invertibile, e $f^{-1}$ è strettamente monotòna dello stesso tipo
(\cref{teo:monotona-invertibile}); senza ipotesi di continuità il viceversa è falso.
\item \textbf{Valore assoluto.} Per $r\ge0$ si ha $\abs{x}\le r\sse-r\le x\le r$
(\cref{prop:modulo-r}); valgono la disuguaglianza triangolare
$\abs{x_{1}+x_{2}}\le\abs{x_{1}}+\abs{x_{2}}$ (\cref{teo:triangolare}) e la
$\bigl|\abs{x_{1}}-\abs{x_{2}}\bigr|\le\abs{x_{1}-x_{2}}$ (\cref{cor:triangolare-inversa}).
\item \textbf{Potenze, esponenziali, logaritmi.} La potenza $x^{n}$ è strettamente
crescente per $x\ge0$ (\cref{prop:potenza-crescente}); per $a>0$, $a\neq1$,
l'esponenziale $a^{x}$ è strettamente monotòna e il logaritmo ne è l'inversa,
$y=\log_{a}x\sse a^{y}=x$ per $x>0$; il logaritmo trasforma prodotti in somme
(\cref{prop:log-proprieta}), ed è questo il principio del regolo calcolatore.
\item \textbf{Trigonometria.} Gli angoli si misurano in radianti; $\sen x$ e $\cos x$
hanno periodo $2\pi$, $\tg x$ ha periodo $\pi$; $\sen^{2}x+\cos^{2}x=1$ e, per
$0<x<\pi/2$, $0<\sen x<x<\tg x$ \eqref{eq:sen-x-tg}, l'unico fatto geometrico che
useremo per derivare le funzioni trigonometriche.
\item \textbf{Il principio di induzione.} Se $\mathcal{P}(1)$ è vera e
$\mathcal{P}(n)\implica\mathcal{P}(n+1)$ per ogni $n\in\N$, allora $\mathcal{P}(n)$ è vera
per ogni $n\in\N$ (\cref{teo:induzione}); con esso si provano le formule delle somme
(\cref{prop:gauss}, \cref{prop:geometrica}, \cref{prop:quadrati},
\cref{prop:telescopica}).
\item \textbf{Disuguaglianza di Bernoulli.} Per $x\ge-1$ e $n\in\N$ si ha
$(1+x)^{n}\ge1+nx$ (\cref{teo:bernoulli}); l'ipotesi $x\ge-1$ è essenziale: per $n=3$ e
$x=-4$ la disuguaglianza è falsa.
\item \textbf{Errori tipici.} $f^{-1}$ non è il reciproco $1/f$; la controimmagine
$f^{-1}(Y)$ ha senso anche se $f$ non è invertibile; la relazione $y^{2}=x$ non definisce
una funzione di $x$; un'affermazione verificata in mille casi non è dimostrata.
\item \textbf{Il calcolatore e i numeri reali.} I numeri in virgola mobile sono un insieme
finito che non soddisfa gli assiomi (\texttt{0.1 + 0.2 == 0.3} è falso): si confrontano
con una tolleranza (\texttt{math.isclose}), e i grafici «per punti» vanno interpretati
alla luce della definizione.
\item \textbf{Per l'ingegnere.} Il radiante è l'unità naturale delle velocità angolari
($v=\omega r$); la bisezione approssima $\sqrt2$ con un errore controllato a priori, che
si dimezza a ogni passo, ed è il metodo che ritroveremo per gli zeri delle funzioni
continue.
\end{puntichiave}

\paragraph{Il filo del discorso.} Il capitolo lascia aperte tre questioni, tutte legate
alla completezza: l'esistenza della radice $n$-esima di ogni numero positivo, finora
soltanto annunciata; la proprietà di Archimede, aggirata nella dimostrazione del
\cref{teo:Q-non-completo}; il significato di $a^{x}$ per $x$ irrazionale. Il prossimo
capitolo trasforma l'assioma di completezza in uno strumento di lavoro, l'estremo
superiore, primo anello della catena dei teoremi di esistenza
(\cref{fig:schema-teoremi}), e con esso risolve tutte e tre le questioni.
